% Chapter 6: Multimodal Knowledge Graph & Vision-RAG

\begin{frame}{Knowledge Graph Architecture}{Conformal-Driven Lazy Execution of the MMKG}
    \begin{block}{Slide Topic \& Explanatory Focus}
        This slide introduces Pipeline 2 (the continuous, data-centric Multimodal Knowledge Graph) and its runtime invocation policy: Conformal-Driven Lazy Execution. Querying a massive graph database on every routine inference call would destroy API throughput and inflate latency. Instead, unambiguous singleton predictions ($|\hat{C}(X)| = 1$) resolve on the Fast Hot-Path using INT8 ONNX models in $\le 22.0$ ms. The MMKG and Vision-RAG pipeline are lazy-loaded exclusively when the Split Conformal Safety Gate detects epistemic uncertainty: severe classification ambiguity ($|\hat{C}(X)| > 3$) where sibling taxa cannot be separated within the $95\%$ confidence bound, or complete novelty/OOD remnants ($|\hat{C}(X)| = 0$).
    \end{block}

    \vspace{0.4cm}
    \statustarget

    \script{Our Multimodal Knowledge Graph is powerful, but querying it takes time. In production, we do not query the graph for every single simple observation. If our conformal safety gate outputs a confident singleton prediction ($|\hat{C}(X)| = 1$), we return the INT8 ONNX prediction on the Fast Hot-Path in under 22 milliseconds. We lazy-load the MMKG only when the conformal gate signals uncertainty—either severe ambiguity ($|\hat{C}(X)| > 3$), or an empty prediction set ($|\hat{C}(X)| = 0$) indicating an anomalous, out-of-distribution specimen.}
    \qa{Why not query the knowledge graph in parallel on every request?}{Parallel querying would saturate database connection pools and waste CPU/GPU cycles on obvious, clear-cut classifications, degrading overall cluster throughput.}
\end{frame}

\begin{frame}{MMKG Pipeline Architecture}{Ingestion \& Retrieval Data Flow}
    \begin{center}
        \resizebox{0.95\textwidth}{!}{
            \begin{tikzpicture}[
                node distance=0.8cm and 1.2cm,
                box/.style={rectangle, draw, rounded corners, align=center, thick, fill=gray!10, minimum height=1.2cm, minimum width=2.5cm},
                arrow/.style={-stealth, thick},
                font=\small\sffamily
            ]
            \providecolor{TaxonPrimary}{RGB}{41, 128, 185}
            \providecolor{TaxonSecondary}{RGB}{39, 174, 96}
            \providecolor{TaxonDark}{RGB}{44, 62, 80}
            \providecolor{TaxonAccent}{RGB}{231, 76, 60}

            \node[box, fill=TaxonPrimary!20, draw=TaxonPrimary] (ingest) {\textbf{Raw Ingestion}};
            \node[box, fill=TaxonSecondary!20, draw=TaxonSecondary, right=of ingest] (filter) {\textbf{Quality Filter} \\ \scriptsize Research Grade, CC licenses, >1024px};
            \node[box, fill=TaxonSecondary!20, draw=TaxonSecondary, right=of filter] (chunk) {\textbf{Visual Chunking} \\ \scriptsize ($I_{\text{child}}$ extraction)};

            \node[box, fill=TaxonDark!20, draw=TaxonDark, below=of chunk, yshift=-0.5cm] (hnsw) {\textbf{Single-Stage Filtered HNSW} \\ \scriptsize Spatial/Seasonal Predicates};
            \node[box, fill=TaxonDark!20, draw=TaxonDark, left=of hnsw] (rrf) {\textbf{RRF Re-ranking} \\ \scriptsize Visual + Geo};
            \node[box, fill=TaxonAccent!20, draw=TaxonAccent, left=of rrf] (vlm) {\textbf{VLM Diagnosis Card} \\ \scriptsize JSON generation};

            \draw[arrow] (ingest) -- (filter);
            \draw[arrow] (filter) -- (chunk);
            \draw[arrow] (chunk) -- (hnsw);
            \draw[arrow] (hnsw) -- (rrf);
            \draw[arrow] (rrf) -- (vlm);
            \end{tikzpicture}
        }
    \end{center}

    \vspace{0.4cm}
    \statustarget

    \script{This diagram maps the full data flow of the MMKG pipeline. Raw data flows through strict quality filters ensuring research-grade, open-license images over 1024 pixels. Then we extract the child organism. The isolated vector is searched in our single-stage filtered HNSW index. Retrieved candidates are fused using Reciprocal Rank Fusion, and finally passed to our local VLM to emit a structured diagnosis card.}
\end{frame}

\begin{frame}{Visual Pre-Processing}{Parent-Child Visual Chunking ($I_{\text{child}}$ vs $I_{\text{parent}}$)}
    \begin{block}{Slide Topic \& Explanatory Focus}
        This slide details our visual entity chunking strategy designed to eliminate background clutter in field photography. Standard visual retrieval across uncropped citizen-science images introduces heavy noise from surrounding soil, rocks, overcast sky, or human hands holding the specimen. TaxonVision employs an upstream zero-shot detector (YOLO-World or OWL-ViT) to isolate the primary organism ($I_{\text{child}}$), to compute clean BioCLIP-2 morphological embeddings ($z_{\text{child}} \in \mathbb{R}^{768}$) without background foliage/soil noise. Simultaneously, the uncropped parent scene context ($I_{\text{parent}}$) is retained for macro-habitat and environmental verification.
    \end{block}

    \vspace{0.4cm}
    \statustarget

    \script{In nature photography, the actual insect or wildflower often takes up less than 10 percent of the frame. The rest is dirt, grass, or human fingers. If you embed the whole image, your vector captures dirt rather than biology. We use YOLO-World or OWL-ViT to detect the specimen and isolate a clean child bounding box crop for fine-grained BioCLIP-2 embedding, while preserving the parent scene image to verify environmental habitat.}
    \qa{What happens if the zero-shot detector finds multiple organisms in one image?}{Each detected child entity generates an independent bounding box and child embedding, allowing multi-organism observations to be analysed separately.}
\end{frame}

\begin{frame}{Visual Pre-Processing}{Parent-Child Visual Chunking Diagram}
    \begin{center}
        \resizebox{0.9\textwidth}{!}{
            \begin{tikzpicture}[
                node distance=1.5cm,
                box/.style={rectangle, draw, rounded corners, align=center, thick, fill=gray!10, minimum height=1.2cm},
                arrow/.style={-stealth, thick},
                font=\small\sffamily
            ]
            \providecolor{TaxonPrimary}{RGB}{41, 128, 185}
            \providecolor{TaxonSecondary}{RGB}{39, 174, 96}
            \providecolor{TaxonDark}{RGB}{44, 62, 80}
            \providecolor{TaxonAccent}{RGB}{231, 76, 60}

            % Parent Image
            \node[box, fill=gray!10, draw=gray, minimum width=4cm, minimum height=3cm] (parent) {};
            \node[anchor=north, font=\small\sffamily\bfseries] at (parent.north) {Full Parent Scene ($I_{\text{parent}}$)};
            \node[anchor=center, text width=3.5cm, align=center] at (parent.center) {\scriptsize Macro-habitat context, \\ background foliage/soil noise};

            % Child Bounding Box inside Parent (conceptual)
            \draw[thick, dashed, draw=TaxonAccent] ([xshift=-1cm,yshift=-0.5cm]parent.center) rectangle ([xshift=1cm,yshift=1cm]parent.center);

            % Detector node
            \node[box, fill=TaxonPrimary!20, draw=TaxonPrimary, right=2cm of parent] (detector) {\textbf{Zero-Shot Detector} \\ \scriptsize YOLO-World / OWL-ViT};

            % Child Crop
            \node[box, fill=TaxonSecondary!20, draw=TaxonSecondary, right=2cm of detector, minimum width=2.5cm, minimum height=2.5cm] (child) {};
            \node[anchor=north, font=\small\sffamily\bfseries] at (child.north) {Child Crop ($I_{\text{child}}$)};
            \node[anchor=center, text width=2cm, align=center] at (child.center) {\scriptsize Clean Organism};

            % BioCLIP-2 Vector
            \node[box, fill=TaxonDark!20, draw=TaxonDark, below=1cm of child] (vector) {\textbf{BioCLIP-2} \\ \scriptsize $z_{\text{child}} \in \mathbb{R}^{768}$};

            \draw[arrow] (parent.east) -- (detector.west);
            \draw[arrow] (detector.east) -- (child.west);
            \draw[arrow] (child.south) -- (vector.north);

            \end{tikzpicture}
        }
    \end{center}

    \vspace{0.4cm}
    \statustarget

    \script{This diagram visually summarises the Parent-Child Visual Chunking process. We take the full parent scene containing background noise, and pass it through a zero-shot detector like YOLO-World or OWL-ViT. This extracts the child bounding box crop, separating the organism from environmental noise, which is then fed to BioCLIP-2 to compute our clean morphological embedding vector of size 768.}
\end{frame}

\begin{frame}{Biogeographical Filtering}{Single-Stage Filtered HNSW with Spatial-Temporal Predicates}
    \begin{block}{Slide Topic \& Explanatory Focus}
        This slide addresses the critical biological trap of convergent evolution—where completely unrelated species on different continents evolve near-identical morphologies (e.g., European wasp vs.\ American cicada killer). Pure visual nearest-neighbour search frequently retrieves geographically impossible candidates. TaxonVision implements single-stage filtered HNSW traversal: spatial constraints (Haversine geodesic radius $R_{\text{max}} \le 500$ km) and temporal constraints ($\Delta_{\text{season}} = 1$ month) are evaluated as boolean predicates natively inside the C++ HNSW graph walk, pruning impossible species before nearest-neighbour calculation completes.
    \end{block}

    \vspace{0.4cm}
    \statustarget

    \script{Nature is full of convergent evolution: two insects on opposite sides of the globe may look identical even though they diverged 50 million years ago. If you only look at pixels, you will misclassify a European wasp as an American cicada killer. In our HNSW vector graph, we evaluate GPS coordinates and seasonal observation month directly inside the graph traversal. If an exemplar is outside a 500-kilometre radius or wrong season, it is pruned instantly before similarity scores are computed.}
    \qa{Why not perform vector search first and filter out GPS locations afterward?}{Post-filtering suffers from candidate starvation: if your top-50 visual neighbours are all from the wrong continent, post-filtering discards all of them and leaves you with zero valid results. Native single-stage filtering guarantees high recall among geographically valid candidates.}
\end{frame}

\begin{frame}{Ensemble Ranking}{Multi-Modal Re-Ranking via Reciprocal Rank Fusion (RRF)}
    \begin{block}{Slide Topic \& Explanatory Focus}
        This slide explores how TaxonVision synthesises visual similarity and biological occurrence priors through Reciprocal Rank Fusion (RRF). After filtered HNSW retrieval yields candidate taxa, the system calculates a unified ranking score combining dense visual cosine similarity ($w_{\text{vis}} = 0.65$) with historical GBIF spatial-temporal occurrence density ($w_{\text{geo}} = 0.35$):
        \[
            \text{RRF}(d) = \frac{w_{\text{vis}}}{k + r_{\text{vis}}(d)} + \frac{w_{\text{geo}}}{k + r_{\text{geo}}(d)}
        \]
        This non-linear rank combination balances fine-grained visual evidence against ecological likelihood without requiring complex score normalisation between divergent probability distributions.
    \end{block}

    \vspace{0.4cm}
    \statustarget

    \script{Once we retrieve candidate species, how do we combine visual similarity with ecological reality? We use Reciprocal Rank Fusion. We assign a 65 percent weight to visual similarity and a 35 percent weight to historical GBIF occurrence density. Even if an exotic species looks slightly closer visually, a high local occurrence density for the native sibling species correctly tilts the ranking towards the true ecological observation.}
    \qa{What is the parameter $k$ in the RRF formula?}{The constant $k$ (typically set to 60) dampens the impact of extreme outliers at top ranks, preventing a single high rank from disproportionately dominating the final ensemble score.}
\end{frame}

\begin{frame}{Clinical-Grade Diagnosis}{VLM Context Synthesis \& Differential Diagnosis Cards}
    \begin{block}{Slide Topic \& Explanatory Focus}
        This slide details our multimodal reasoning interface for human taxonomists. When an ambiguous observation reaches triage, the query crop ($I_{\text{child}}$) and the top-3 re-ranked exemplar specimens are injected into a local Vision-Language Model (e.g., Qwen2-VL-2B). Constrained by formal grammar-based sampling, the VLM generates a structured JSON differential diagnosis card detailing distinguishing anatomical discriminators, ecological likelihoods, and CAM activation overlays. These diagnosis cards are rendered directly into Label Studio, transforming triage into rapid, high-precision expert verification.
    \end{block}

    \vspace{0.4cm}
    \statustarget

    \script{When a human taxonomist is asked to review an ambiguous specimen, we don't just dump raw images on their screen. We pass the query crop and the top-3 reference exemplars into a local Vision-Language Model like Qwen2-VL-2B. Constrained by grammar-based decoding, the VLM generates a structured JSON differential diagnosis card: it details specific anatomical discriminators, shows CAM overlays, and allows the expert to confirm the species in seconds.}
    \qa{Why use grammar-based sampling with the Vision-Language Model?}{Unconstrained VLMs hallucinate text and output unpredictable schemas. Grammar-based sampling forces the model to emit strictly valid JSON conforming to our differential diagnosis schema.}
\end{frame}

\begin{frame}{Taxonomic Agility}{Zero-Retraining Taxonomic Mutations via DuckDB}
    \begin{block}{Slide Topic \& Explanatory Focus}
        This slide examines how TaxonVision handles dynamic changes in biological taxonomy without expensive GPU model retraining. In biological science, taxonomic committees frequently alter classifications: merging species (taxonomic lumps) or dividing a single species into multiple distinct species (taxonomic splits). Because visual vectors reside in HNSW while taxonomic relationships and observation metadata reside in our decoupled DuckDB and Parquet architecture, taxonomic lumps execute as $\mathcal{O}(1)$ relational updates with zero GPU re-embedding cost. Taxonomic splits mark historical nodes with \texttt{requires\_reannotation = true}, sending them to triage before re-indexing.
    \end{block}

    \vspace{0.4cm}
    \statustarget

    \script{In real biology, taxonomy is not static. Scientists frequently merge two species into one, or split one species into three. In a standard deep learning model, every split or merge requires retraining the entire neural network from scratch. Because our architecture decouples HNSW vector storage from DuckDB relational metadata, a taxonomic merge is an instant O(1) database update. Our vectors remain untouched with zero GPU compute cost.}
    \qa{What happens during a taxonomic split when one species becomes two?}{Historical observation nodes for that taxon in DuckDB are flagged with \texttt{requires\_reannotation = true} and purged from active traversal. Their crops are routed to Label Studio for expert re-sorting, after which they are re-registered under their new respective taxon keys.}
\end{frame}
